\section*{3. 同余公式: 陈述及初等的等价形式}

感谢 Grothendieck 和 Deligne 向我解释 (3.1.1) 的变体
(3.1.2)、(3.2.1)--(3.2.2) 及 (3.3.1).

\subsection*{3.0.}
设 $X$ 是 $\mathbb F_q$ 上的有限型概形.
闭点 $y\in X$ 的次数定义为其剩余域扩张 $\mathbb F_q(y)/\mathbb F_q$ 的次数.
$X/\mathbb F_q$ 的 zeta 函数是 $\mathbb Z[[t]]$ 的元素,
由下式定义 (亦参见 SGA~5, XIII~3.1)

\begin{equation}\tag{3.0.1}
Z(t,X/\mathbb F_q)=\prod_y(1-t^{\deg(y)})^{-1}.
\end{equation}

“模 $p$ 的 zeta 函数”始终指 $Z(t,X/\mathbb F_q)$ 在“模 $p$ 约化”映射下
于 $\mathbb F_q[[t]]$ 中的像

\[
\mathbb Z[[t]]\longrightarrow\mathbb F_q[[t]].
\]

\subsubsection*{定理 3.1.}
设 $X$ 是 $\mathbb F_q$ 上的固有概形. 则

\begin{equation}\tag{3.1.1}
Z(t,X/\mathbb F_q)
\equiv
\prod_{i\geq0}
\det\!\left(1-t\mathfrak F_q\mid H^i(X,\mathcal O_X)\right)^{(-1)^{i+1}}
\pmod p,
\end{equation}

或者, 由命题~2.2.4 等价地,

\begin{equation}\tag{3.1.2}
Z(t,X/\mathbb F_q)
\equiv
\prod_{i\geq0}
\det\!\left(1-t\varphi_q^{-1}\mid
H^i_{\mathrm{et}}(\bar X,\mathbb Z/p\mathbb Z)\right)^{(-1)^{i+1}}
\pmod p.
\end{equation}

\subsubsection*{推论 3.2.}
设 $X$ 是 $\mathbb F_q$ 上的固有概形. 则

\begin{equation}\tag{3.2.1}
\sum_{i\geq0}(-1)^i
\operatorname{Tr}\!\left(\mathfrak F_q\mid H^i(X,\mathcal O_X)\right)
\equiv \#X(\mathbb F_q)\pmod p,
\end{equation}

或者, 由命题~2.2.4 等价地,

\begin{equation}\tag{3.2.2}
\sum_{i\geq0}(-1)^i
\operatorname{Tr}\!\left(\varphi_q^{-1}\mid
H^i_{\mathrm{et}}(\bar X,\mathbb Z/p\mathbb Z)\right)
\equiv \#X(\mathbb F_q)\pmod p.
\end{equation}

事实上, 在同余式 (3.1.1) 及 (3.1.2) 中取 $t$ 的系数即可.

\subsubsection*{定理 3.3.}
设 $X$ 是 $\mathbb F_q$ 上的固有概形, 用 $H^i_{\mathrm{DR}}(X/\mathbb F_q)$
表示 $\mathbb F_q$ 向量空间 $H^i(X,\Omega^\bullet_{X/\mathbb F_q})$,
其上有 $\mathfrak F_q$ 作用. 则

\begin{equation}\tag{3.3.1}
Z(t,X/\mathbb F_q)
\equiv
\prod_{i\geq0}
\det\!\left(1-t\mathfrak F_q\mid
H^i_{\mathrm{DR}}(X/\mathbb F_q)\right)^{(-1)^{i+1}}
\pmod p.
\end{equation}

事实上, 考虑谱序列

\begin{equation}\tag{3.3.2}
E_1^{p,q}=H^q(X,\Omega^p_{X/\mathbb F_q})
\Longrightarrow H^{p+q}_{\mathrm{DR}}(X/\mathbb F_q).
\end{equation}

算子 $\mathfrak F_q$ 通过 $\mathrm{Fr}_q^X:X\longrightarrow X$ 诱导的自同态
$\mathfrak F_q=(\mathrm{Fr}_q^X)^*$ 作用于其上,
并在 $p\ne0$ 时将 $E_1^{p,q}$ 化为零. 因而 (3.3.1) 与 (3.1.1) 等价.

再作一点补充. 由于取“半单部分” (参见 1.0.7) 是正合函子,
(3.3.2) 给出谱序列

\begin{equation}\tag{3.3.3}
(E_1^{p,q})_{\mathrm{ss}}
\Longrightarrow H^{p+q}_{\mathrm{DR}}(X/\mathbb F_q)_{\mathrm{ss}}.
\end{equation}

由于 $p\ne0$ 时 $(E_1^{p,q})_{\mathrm{ss}}=0$, 对每个 $i$ 得到同构

\begin{equation}\tag{3.3.4}
H^i_{\mathrm{DR}}(X/\mathbb F_q)_{\mathrm{ss}}
\xrightarrow{\ \sim\ }
H^i(X,\mathcal O_X)_{\mathrm{ss}}
\end{equation}

它与 $\mathfrak F_q$ 可交换. 因此

\begin{equation}\tag{3.3.5}
\det\!\left(1-t\mathfrak F_q\mid H^i_{\mathrm{DR}}(X/\mathbb F_q)\right)
=
\det\!\left(1-t\mathfrak F_q\mid H^i(X,\mathcal O_X)\right).
\end{equation}
